% Part: first-order-logic
% Chapter: natural-deduction
% Section: proof-theoretic-notions

\documentclass[../../../include/open-logic-section]{subfiles}

\begin{document}

\iftag{FOL}
      {\olfileid{fol}{ntd}{ptn}}
      {\olfileid{pl}{ntd}{ptn}}

\olsection{Bewijstheoretische begrippen}

\begin{editorial}
In deze paragraaf worden de definities van de afleidbaarheidsrelatie
en van consistentie voor natuurlijke deductie bijeengebracht.
\end{editorial}

\begin{explain}
Zoals we een aantal belangrijke semantische begrippen hebben
gedefinieerd (geldigheid, gevolg en vervulbaarheid), definiëren we nu
de overeenkomstige \emph{bewijstheoretische begrippen}. Deze worden
niet gedefinieerd aan de hand van het waar zijn van !!{sentence}s in
!!{structure}s, maar aan de hand van de !!{derivability} of
!!{nonderivability} van bepaalde !!{sentence}s uit andere. De
ontdekking dat deze begrippen samenvallen, was van groot belang. Dat
zij inderdaad samenvallen, is de inhoud van de
\emph{correctheidsstelling} en de \emph{volledigheidsstelling}.
\end{explain}


\begin{defn}[Stellingen]
!!^a{sentence}~$!A$ is een \emph{stelling} als er in de natuurlijke
deductie !!a{derivation} van~$!A$ bestaat waarin alle aannamen zijn
!!{discharged}. We schrijven $\Proves !A$ als $!A$ een stelling is en
$\Proves/ !A$ als dat niet zo is.
\end{defn}

\begin{defn}[!!^{derivability}]
!!^a{sentence} $!A$ is \emph{!!{derivable} uit} een verzameling
!!{sentence}s~$\Gamma$, genoteerd als $\Gamma \Proves !A$, als er
!!a{derivation} met conclusie~$!A$ bestaat waarin elke aanname ofwel
!!{discharged} is, ofwel tot~$\Gamma$ behoort. Als $!A$ niet
!!{derivable} is uit $\Gamma$, schrijven we $\Gamma \Proves/ !A$.
\end{defn}

\begin{defn}[Consistentie]
Een verzameling !!{sentence}s~$\Gamma$ is \emph{inconsistent} dan en
slechts dan als $\Gamma \Proves \lfalse$. Als $\Gamma$ niet
inconsistent is, dat wil zeggen als $\Gamma \Proves/ \lfalse$, noemen
we haar \emph{consistent}.
\end{defn}

\begin{prop}[Reflexiviteit]
\ollabel{prop:reflexivity}
Als $!A \in \Gamma$, dan $\Gamma \Proves !A$.
\end{prop}

\begin{proof}
De aanname $!A$ op zichzelf is !!a{derivation} van~$!A$ waarin elke
!!{undischarged} aanname (namelijk $!A$) tot~$\Gamma$ behoort.
\end{proof}
  

\begin{prop}[Monotonie]
\ollabel{prop:monotonicity}
Als $\Gamma \subseteq \Delta$ en $\Gamma \Proves !A$, dan $\Delta
\Proves !A$.
\end{prop}

\begin{proof}
Elke !!{derivation} van $!A$ uit $\Gamma$ is tevens !!a{derivation}
van $!A$ uit~$\Delta$.
\end{proof}

\begin{prop}[Transitiviteit]
\ollabel{prop:transitivity}
Als $\Gamma \Proves !A$ en $\{!A\} \cup \Delta \Proves
!B$, dan $\Gamma \cup \Delta \Proves !B$.
\end{prop}

\begin{proof}
Als $\Gamma \Proves !A$, bestaat er !!a{derivation}~$\delta_0$
van~$!A$ waarvan alle !!{undischarged} aannamen in~$\Gamma$ liggen.
Als $\{!A\} \cup \Delta \Proves !B$, bestaat er !!a{derivation}
~$\delta_1$ van~$!B$ waarvan alle !!{undischarged} aannamen in
~$\{!A\} \cup \Delta$ liggen. Beschouw nu:
\begin{prooftree}
  \AxiomC{$\Delta, \Discharge{!A}{1}$}
  \RightLabel{$\delta_1$}
  \DeduceC{$!B$}
  \DischargeRule{\Intro{\lif}}{1}
  \UnaryInfC{$!A \lif !B$}
  \AxiomC{$\Gamma$}
  \RightLabel{$\delta_0$}
  \DeduceC{$!A$}
  \RightLabel{\Elim{\lif}}
  \BinaryInfC{$!B$}
\end{prooftree}
De !!{undischarged} aannamen behoren nu alle tot $\Gamma \cup
\Delta$, en dit toont dus $\Gamma \cup \Delta \Proves !B$ aan.
\end{proof}

Als $\Gamma = \{!A_1, !A_2, \ldots, !A_k\}$ een eindige verzameling
is, mogen we de verkorte notatie
$!A_1,!A_2,\ldots,!A_k \Proves !B$ gebruiken voor
$\Gamma \Proves !B$. In het bijzonder
betekent $!A \Proves !B$ dat $\{!A\} \Proves !B$.

Merk op dat uit $\Gamma \Proves !A$ en $!A
\Proves !B$ volgt dat $\Gamma \Proves !B$. Hieruit volgt ook dat, als
$!A_1, \dots, !A_n \Proves !B$ en $\Gamma \Proves !A_i$ voor
elke~$i$, dan $\Gamma \Proves !B$.

\begin{prop}
\ollabel{prop:incons}
De volgende uitspraken zijn equivalent.
    \begin{enumerate}
      \item \( \Gamma \) is inconsistent.
      \item \( \Gamma \Proves {!A} \) voor elke !!{sentence}~\( {!A} \).
      \item Zowel \( \Gamma \Proves {!A} \) als
        \( \Gamma \Proves \lnot {!A} \) geldt voor een zekere
        !!{sentence}~\( {!A} \).
    \end{enumerate}
\end{prop}

\begin{proof}
Oefening.
\end{proof}

\tagprob{FOL}
\begin{prob}
Bewijs \olref[fol][ntd][ptn]{prop:incons}
\end{prob}
\tagendprob

\tagprob{notFOL}
\begin{prob}
Bewijs \olref[pl][ntd][ptn]{prop:incons}
\end{prob}
\tagendprob

\begin{prop}[Compactheid]
  \ollabel{prop:proves-compact}
  \begin{enumerate}
  \item Als $\Gamma \Proves !A$, dan bestaat er een eindige
    deelverzameling $\Gamma_0 \subseteq \Gamma$ waarvoor
    $\Gamma_0 \Proves !A$.
  \item Als elke eindige deelverzameling van~$\Gamma$ consistent is,
    dan is $\Gamma$ consistent.
  \end{enumerate}
\end{prop}

\begin{proof}
  \begin{enumerate}
    \item Als $\Gamma \Proves !A$, bestaat er !!a{derivation}
      ~$\delta$ van~$!A$ uit~$\Gamma$. Laat $\Gamma_0$ de verzameling
      !!{undischarged} aannamen van~$\delta$ zijn. Omdat elke
      !!{derivation} eindig is, kan $\Gamma_0$ slechts eindig veel
      !!{sentence}s bevatten. Dus is $\delta$ !!a{derivation} van
      ~$!A$ uit een eindige~$\Gamma_0 \subseteq \Gamma$.
    \item Dit is de contrapositie van (1) voor het bijzondere geval
      $!A \ident \lfalse$.
  \end{enumerate}
\end{proof}
\end{document}
